\subsection{Exact branch seeds and continuous synthetic truth}
Supplying the exact earliest branch separates initialization error from
path representation and local optimization. Four identified Loma rays,
diagnostic rows 5, 49, 57, and 88, have continuous first arrivals along the
critical guide at $x=0$. Circular arcs and the straight guide are sampled
at five spacings from 0.5 to 0.03125 km. One calculation uses the unchanged
tracer, including its initial uniform resampling. A second copied tracer
retains the supplied nodes initially; all interior nodes remain free during
optimization. The action, derivatives, line search, and 1,000-update budget
remain unchanged.

Only one of twenty optimized uniform-seed paths meets the local stopping
rule; none of twenty retained-seed paths does. At 0.03125 km, all four
uniform-seed results nevertheless have time errors below 4.750 microseconds
and source-vector errors below $4.35\times10^{-7}$ s/km. The finest retained
row-5 path accepts no update, despite time error 4.696 microseconds and
source-vector error about $2\times10^{-15}$ s/km. Its ordinary gradient norm
is 0.0751 s/km, and collective normal displacements of its interface nodes
have unequal one-sided action slopes. The correct arrival branch therefore
does not resolve the nonsmooth optimization problem. These controls also
show that initial uniform resampling can damage a coarse guide seed; its
additional row-5 time error falls from 473.730 microseconds at 0.5 km to
0.05368 microseconds at 0.03125 km. Exact polygon integration agrees with
production action within $2.25\times10^{-13}$ s. All paths, original flags,
identities, and isolated source changes are retained in
\path{data/development/loma_geometry_preflight/reference_seeds_001/}.

Continuous layered enumeration was then extended to all 37,648 pairs in
the prepared 724-source, 52-receiver geometry. The first attempt completed
every candidate enumeration, but eleven direct candidates failed the
unchanged $10^{-8}$ km distance-closure requirement. Their first-arrival
values remain blank. All eleven use source 357, event
\texttt{NC/20091139}, at $x=-0.011677621261526383$ km, just inside the
6 km/s constant-speed region. Solving for $p v_{\max}$ close to one loses
the small critical offset in binary64 arithmetic. A separate correction
solves for $\delta=1/p-v_{\max}\geq0$ and forms the phase speed by exact
rational addition before evaluating the segment integrals. Twenty-one
isolated cases, including all eleven failed geometries and ten invented
near-interface controls, close within $7.11\times10^{-15}$ km. Independent
80-digit integrals agree in time within $3.91\times10^{-15}$ s and in
source-vector norm within $2.87\times10^{-17}$ s/km. The failed first attempt
is preserved under \path{data/development/loma_layered_full_001/}.

The corrected second enumeration passes the full geometry and independent
native-identity review: all 37,648 references are valid, with no near ties.
Their first arrivals comprise 25,724 direct paths, 5,951 turning paths,
5,474 critical guides, and 499 direct paths in constant-speed regions.
The maximum closure residual across all 66,302 candidates is
$3.13\times10^{-9}$ km. The eleven formerly invalid references are the only
validity changes. Previously valid branch kinds remain unchanged; their
largest time change is 1.55 ns. No source or receiver is removed. Both
compressed identified streams, the executed reference source, and the
separate completion, physical-consistency, and independent identity reviews
are retained in \path{data/development/loma_layered_full_002/}.

An explicit analytic-truth input mode permits a separately reviewed,
complete continuous reference to supply synthetic travel times. The C++
consumer verifies the model, ordered source--receiver support, every native
metadata field, finite values, unique valid references, and physical speed
bounds. It preserves the existing random-draw, quantization, and differential
observation order. Its invented contract suite passes 55 checks, and six
default-mode output files remain byte-identical to the frozen preceding
executable. Analytic mode writes reference status separately and leaves
polygon-ray status fields blank. This changes the declared synthetic truth
calculation; it does not repair or accept the inverse forward operator.

\subsection{Historical waveform metadata and a rejected pilot pair}
The existing Hayward catalog pair at differential row 15108 links events
\texttt{NC/238035} and \texttt{NC/30083101} at
\texttt{NC.JPR.02.EHZ}. EarthScope routing identifies NCEDC as the archive
for the two requested windows. The preserved miniSEED records each contain
500 contiguous samples at 100 Hz over a half-open five-second interval.
Full response metadata cover the distinct 1991 and 1995 epochs; their
three-stage response contents agree. Event, origin, pick, association,
channel, and epoch identities accompany the raw files and derived records.
Missing provider response resource IDs remain explicitly missing.

The written precursor method uses a tapered 2.56 s window and requires
squared coherence above 0.9 at at least half the frequency bins between
2 and 10 Hz (Waldhauser and Ellsworth, 2002, section 2.2, paragraph 16).
For the primary pilot calculation, the unspecified choices are declared:
a window beginning 0.5 s before the catalog P pick, demeaning, a Hann taper,
and five-bin spectral averaging. No additional filter, response removal,
or resampling is applied. Only two of twenty frequency bins pass, where ten
are required. All nine declared window/smoothing variants fail, with zero
to four passing bins. Consequently, no waveform difference is adopted.
A line through the two coherent phase bins has two free parameters; its
near-zero residual does not measure timing precision. The source, spectra,
controls, and failed result are in
\path{data/development/waveform_correlation_pilot_001/}.

An offline audit retains all 14,262 picks and 44,718 catalog differences.
Historical operation records identify \texttt{NC.CMW1}, but do not provide
calibration: the saved EarthScope and NCEDC StationXML responses start in
2000, whereas 117 retained picks and 199 differences involve its 1995--1998
epochs. No later response is backdated. Among 398 picks whose channel names
do not end in Z, provider orientation metadata identify 227 as vertical,
54 as nonvertical, and 117 as unknown. Channel names alone therefore cannot
define the vertical waveform selection. A further 89 picks and 542
differences precede the archive's documented March 1984 start; this date
is a coverage lead, not proof that all earlier data are absent. These audit
categories remove no catalog observations and establish no study exclusion.
Further retrieval is paused for clarification under the user's metadata
rule. Full native rows and unsubmitted requests are saved in
\path{data/authoritative/waveforms/cohort_plan_001/}.
